\chapter{Time-Series Databases and the Alternatives}\label{ch:pl:time-series-databases-and-the-alternatives}

Five engines answered the same five questions on the same month of quotes, and each answer was checked against the other four
before anything was timed. The row store with an index found one symbol's last quote before noon in less than a tenth of a
millisecond, nearly eighty times faster than the columnar engine, and took over three hundred milliseconds to average the
spread of every symbol over the month, over forty times slower than the columnar engine. Neither result is a property of the product; both are
properties of the question. This chapter looks at what the databases sold for time series promise, sorts market-data queries
into a few shapes, explains why an engine is fast on some shapes and slow on others, and builds the benchmark
(\texttt{firm.tsbench}) that a firm should run on its own data before it chooses.

\section{What a time-series database promises}

\begin{definition}[Time-series database]\label{def:pl:time-series-databases-and-the-alternatives:tsdb}
A \emph{time-series database}\index{time-series database} is a database designed for records keyed by time that mostly arrive
in time order and are mostly read by time ranges: it stores them partitioned and ordered by time, appends quickly, and offers
time-specific operations (aggregation into time intervals, as-of joins, retention by age).
\end{definition}

The chapter-4 \omterm{def:pl:a-tick-store-on-open-formats:store}{tick store} has every one of those properties and is not a database product: it is files, an embedded engine and
a few conventions. That is the first thing to understand about the category. What a \omterm{def:pl:time-series-databases-and-the-alternatives:tsdb}{time-series database} adds is packaging ---
a server that ingests, compacts, retains and answers queries without the firm writing the pieces --- and, in the better ones,
engineering of the same ideas: columnar storage, time partitions, \omterm{def:pl:columnar-formats:pushdown}{zone maps}, vectorised execution. The operations it names are
the ones every market-data store needs.

\begin{definition}[Time bucket, downsampling]\label{def:pl:time-series-databases-and-the-alternatives:bucket}
A \emph{time bucket}\index{time bucket} is one of the equal, aligned intervals into which a time axis is cut (the minutes of a
day); bucketing assigns each record to its interval. \emph{Downsampling}\index{downsampling} replaces the records of each bucket
by a summary (the last value, the count, an open-high-low-close bar), storing or returning one row per bucket.
\end{definition}

One-minute bars are the standard example: One Quant Book~7's time bars (chapter~2) are downsampled ticks. A store that keeps
tick history for months and minute bars for years downsamples as part of its \omterm{def:pl:capturing-and-storing-tick-data:tier}{retention policy} (chapter~2): the raw ticks age
out, the buckets stay.

\section{Query shapes of market data}

Almost every query a research or trading platform sends to its market data is one of a handful of shapes, and each shape asks
something different of the storage (\cref{fig:pl:time-series-databases-and-the-alternatives:shapes}).

\begin{omfigure}
\begin{tikzpicture}[font=\scriptsize,>=Stealth]
\draw[->] (0,0) -- (10.6,0) node[below left]{time (a month)};
\draw[->] (0,0) -- (0,4.1) node[left,align=right]{symbols};
\fill[gray!10] (0.05,0.1) rectangle (10.25,3.8);
\draw[gray!60,dashed] (0.05,0.1) rectangle (10.25,3.8);
\foreach \y in {0.5,1.0,...,3.5} {\draw[gray!25] (0,\y) -- (10.3,\y);}
\fill[omThm!70] (5.02,1.95) circle (2.2pt);
\node[omThm,anchor=south east] at (4.95,2.08) {point};
\fill[omDef!40] (5.3,1.9) rectangle (5.9,2.1);
\node[omDef,anchor=south west] at (5.35,2.12) {range};
\fill[omMeth!40] (6.5,1.4) rectangle (7.2,1.6);
\foreach \x in {6.6,6.7,6.8,6.9,7.0,7.1} {\draw[omMeth] (\x,1.4) -- (\x,1.6);}
\node[omMeth,anchor=north] at (6.85,1.38) {bars (buckets)};
\fill[omProp!35] (3.2,0.2) rectangle (3.6,3.7);
\node[omProp,anchor=south] at (3.4,3.72) {as-of join (one day, all symbols)};
\node[anchor=south east,fill=gray!10,inner sep=1pt] at (10.2,0.15) {cross-section: every row (dashed box)};
\end{tikzpicture}

\omcaption{Five query shapes on a month of quotes (symbols by time): a point lookup (the last quote of one symbol before one
time), a range scan (one symbol, one hour), \omterm{def:pl:time-series-databases-and-the-alternatives:bucket}{time buckets} (one symbol, one day, one row a minute), an as-of join (every trade of a
day with its symbol's quote), and a cross-section (every row of the month, one number per symbol). Each touches a different
region, and a different layout serves each.}\label{fig:pl:time-series-databases-and-the-alternatives:shapes}
\end{omfigure}

\begin{table}[ht]
\centering\small
\begin{tabular}{llll}
\toprule
shape & rows touched & what makes it fast & example \\
\midrule
point & one & an index or a sorted key & last quote before an order \\
range & a contiguous slice & sort order, \omterm{def:pl:columnar-formats:pushdown}{zone maps} & one symbol's hour \\
bars & a slice, summarised & sort order, vectorised aggregation & minute bars for a chart \\
as-of join & two slices, merged & both sorted by time within symbol & trades with quotes \\
cross-section & all & columns, vectorised scans & average spread by symbol \\
\bottomrule
\end{tabular}
\caption{The five query shapes, the rows each touches, and what an engine needs to answer it quickly.}\label{tab:pl:time-series-databases-and-the-alternatives:shapes}
\end{table}

\section{In-process engines against servers}

A database was once a server: a separate process, often on a separate machine, that clients sent queries to over a network.
The analytical engines most research platforms now use are not.

\begin{definition}[Embedded database]\label{def:pl:time-series-databases-and-the-alternatives:embedded}
An \emph{embedded database}\index{embedded database} runs inside the process that queries it, as a library: no server, no
network, no serialisation of results, and data shared with the host program's memory where the formats allow.
\end{definition}

SQLite, the most deployed database there is, is embedded and row-oriented; DuckDB is embedded and columnar; Polars is a dataframe
library with a query optimiser, which makes it an embedded engine too. An embedded engine gives up what a server gives ---
concurrent writers, access control, one copy of the data for many users --- and gains the round trips and the conversions it
does not make (\cref{fig:pl:time-series-databases-and-the-alternatives:embedded}). For a researcher's queries on files in a
\omterm{def:pl:a-tick-store-on-open-formats:store}{tick store}, that is the right trade; for a service many teams write to, it is not (chapter~25).

\begin{omfigure}
\begin{tikzpicture}[font=\scriptsize,>=Stealth,box/.style={draw,rounded corners=2pt,minimum height=0.7cm,inner sep=3pt,align=center}]
\draw[rounded corners=3pt,gray] (-0.3,-0.9) rectangle (4.1,1.0);
\node[font=\scriptsize\itshape,anchor=south west] at (-0.3,1.0) {one process};
\node[box,fill=omDef!12] (a) at (0.9,0) {research code};
\node[box,fill=omMeth!15] (e) at (3.1,0) {engine\\(library)};
\draw[<->] (a) -- (e);
\node[box,fill=gray!10] (f) at (3.1,-1.9) {columnar files};
\draw[->] (f) -- (e);
\draw[rounded corners=3pt,gray] (5.0,-0.9) rectangle (7.0,1.0);
\draw[rounded corners=3pt,gray] (9.6,-0.9) rectangle (12.0,1.0);
\node[font=\scriptsize\itshape,anchor=south west] at (5.0,1.0) {client};
\node[font=\scriptsize\itshape,anchor=south west] at (9.6,1.0) {server};
\node[box,fill=omDef!12] (c) at (6.0,0) {research code};
\node[box,fill=omMeth!15] (s) at (10.8,0) {database\\server};
\draw[->] (c.north east) to[bend left=15] node[above]{query} (s.north west);
\draw[->] (s.south west) to[bend left=15] node[below,align=center]{rows,\\ serialised} (c.south east);
\node[box,fill=gray!10] (d) at (10.8,-1.9) {its own storage};
\draw[->] (d) -- (s);
\end{tikzpicture}

\omcaption{An embedded engine (left) runs inside the program that queries it and reads the files directly; a database server
(right) is another process, reached over a connection, which returns its results serialised. The server can serve many users
and writers; the embedded engine skips the round trip and the conversion.}\label{fig:pl:time-series-databases-and-the-alternatives:embedded}
\end{omfigure}

The other difference between engines is how they execute a plan. A classic engine passes one row at a time from operator to
operator, calling a function per row per operator. A columnar engine passes batches of values of one column.

\begin{definition}[Vectorised query execution, late materialisation]\label{def:pl:time-series-databases-and-the-alternatives:vectorised}
\emph{Vectorised query execution}\index{vectorised query execution} runs each operator of a query plan on a batch of values of a
column at a time (a vector of a few thousand), so that the per-call overhead is paid per batch and the inner loops are tight
loops over arrays. \emph{Late materialisation}\index{late materialisation} delays assembling result rows from columns until
after the filters: operators pass positions of qualifying rows, and only the surviving rows' other columns are ever read.
\end{definition}

Both ideas come from the research line of the MonetDB/X100 engine (2005) and the column stores that followed; they are why a
columnar engine scans a million rows in milliseconds. They are also why it is slower than an index on a point lookup: a
vectorised scan of a sorted file still opens a file, reads a footer and decodes at least a block, where a B-tree walks three
pages to one row.

\begin{omfigure}
\begin{tikzpicture}[font=\scriptsize,>=Stealth,op/.style={draw,rounded corners=2pt,minimum width=1.9cm,minimum height=0.6cm,fill=omDef!10}]
\node[op] (s1) at (0,2.4) {scan};
\node[op] (f1) at (0,1.2) {filter};
\node[op] (g1) at (0,0) {aggregate};
\draw[->] (s1) -- node[right]{one row} (f1);
\draw[->] (f1) -- node[right]{one row} (g1);
\node[font=\scriptsize\itshape] at (0,3.1) {row at a time};
\node[op,fill=omMeth!15] (s2) at (5.0,2.4) {scan};
\node[op,fill=omMeth!15] (f2) at (5.0,1.2) {filter};
\node[op,fill=omMeth!15] (g2) at (5.0,0) {aggregate};
\draw[->] (s2) -- node[right]{2\,048 values of a column} (f2);
\draw[->] (f2) -- node[right]{positions that pass} (g2);
\node[font=\scriptsize\itshape] at (5.0,3.1) {vectorised};
\end{tikzpicture}

\omcaption{Row-at-a-time and vectorised execution of the same plan. The vectorised engine calls each operator once per batch of
a column and passes positions instead of rows (\omterm{def:pl:time-series-databases-and-the-alternatives:vectorised}{late materialisation}), so the other columns of rows that fail the filter are never
read.}\label{fig:pl:time-series-databases-and-the-alternatives:vector}
\end{omfigure}

\begin{proposition}[No engine wins every shape]\label{prop:pl:time-series-databases-and-the-alternatives:shapes}
Let an engine answer a query in time $c_0 + c_r\,r$, where $r$ is the number of rows it must touch and $c_0$ the fixed cost of
starting (opening files, planning, reading metadata). An indexed row store has a small $c_0$ and a large $c_r$; a vectorised
columnar engine has a larger $c_0$ and a much smaller $c_r$. The row store is faster below $r^\star = (c_0^{\mathrm{col}} -
c_0^{\mathrm{row}})/(c_r^{\mathrm{row}} - c_r^{\mathrm{col}})$ rows touched, the columnar engine above.
\end{proposition}

\begin{proof}
Equate the two affine costs and solve for $r$; below the crossing the smaller intercept wins, above it the smaller slope.
\end{proof}

The model is crude --- a real engine's cost depends on layout, selectivity and cache --- but its prediction is what the benchmark
finds: point lookups and short ranges on one side, cross-sections on the other.

\section{A benchmark that means something}

\begin{method}[Benchmarking engines on market data]\label{met:pl:time-series-databases-and-the-alternatives:bench}
(1) Use the firm's own data and query shapes, not a vendor's demonstration. (2) Check that every engine returns the same answer
to every query before timing any of them: a fast wrong answer is the commonest benchmark result. (3) Run the engines interleaved,
every engine once per round, so that a drifting machine affects all alike. (4) Report the first run on a fresh engine (cold)
apart from the median of the rest (warm), and the spread. (5) State the machine, the load and what was in the page cache.
\end{method}

\texttt{firm.tsbench} implements the method for five engines: DuckDB over the month's Parquet files, partitioned by date and
sorted by symbol and time; Polars scanning the same files lazily; SQLite, embedded and row-oriented, with an index on symbol and
time; pandas, eager and in memory; and a hand-written numpy engine on the arrays sorted by symbol and time with a per-symbol
offset index. \Cref{lst:pl:time-series-databases-and-the-alternatives:duck} is one engine's five queries.

\omcode{code/firm/tsbench/firm_tsbench.py}{76}{101}{The five query shapes as DuckDB answers them: a sorted lookup, a range, time
buckets with the last mid of each minute, an as-of join and a cross-sectional average.}\label{lst:pl:time-series-databases-and-the-alternatives:duck}

The data are the chapter-3 month of quotes (a million rows, fifty symbols, twenty days) and 40\,000 trades drawn from it, each
one nanosecond after a quote. All five engines agree on all five answers: 1 row for the point lookup, 153 for the hour of
S007, 359 one-minute bars for its day, 2\,000 trades joined, 50 average spreads.

\begin{omfigure}
\begin{tikzpicture}
\begin{axis}[width=12cm,height=6.4cm,ybar,bar width=5pt,ymode=log,log origin=infty,ymin=0.03,ymax=1000,
  ylabel={median time (ms, log scale)},xtick={0,1,2,3,4},
  xticklabels from table={figdata/platforms/05-time-series-databases-and-the-alternatives/measured_tsbench_wide.csv}{query},
  tick label style={font=\small},label style={font=\small},grid=major,grid style={gray!20},enlarge x limits=0.15,
  legend columns=5,legend style={font=\footnotesize,at={(0.5,-0.18)},anchor=north,draw=none,/tikz/every even column/.append style={column sep=6pt}},
  table/col sep=comma]
\addplot[fill=omDef!55,draw=omDef] table[x=k,y=duckdb]{figdata/platforms/05-time-series-databases-and-the-alternatives/measured_tsbench_wide.csv};
\addplot[fill=omMeth!55,draw=omMeth] table[x=k,y=polars]{figdata/platforms/05-time-series-databases-and-the-alternatives/measured_tsbench_wide.csv};
\addplot[fill=omThm!55,draw=omThm] table[x=k,y=sqlite]{figdata/platforms/05-time-series-databases-and-the-alternatives/measured_tsbench_wide.csv};
\addplot[fill=omProp!55,draw=omProp] table[x=k,y=pandas]{figdata/platforms/05-time-series-databases-and-the-alternatives/measured_tsbench_wide.csv};
\addplot[fill=gray!55,draw=gray] table[x=k,y=numpy]{figdata/platforms/05-time-series-databases-and-the-alternatives/measured_tsbench_wide.csv};
\legend{DuckDB,Polars,SQLite,pandas,numpy}
\end{axis}
\end{tikzpicture}

\omcaption{Median warm time of each query shape on each engine, seven interleaved rounds, answers checked equal first. Measured
on a laptop (Intel Core Ultra 7 155H) under WSL2, one thread, files in the page cache, machine otherwise idle. Data:
\texttt{bench\_tsbench.py}.}\label{fig:pl:time-series-databases-and-the-alternatives:times}
\end{omfigure}

\Cref{fig:pl:time-series-databases-and-the-alternatives:times} is the result, and the proposition's prediction holds. SQLite's
index answers the point lookup in 0.05~ms and the hour in 0.3~ms, DuckDB takes 3.8 and 2.2~ms for the same two, dominated by
opening and planning; on the cross-section the order reverses, SQLite needing 313~ms to read a million rows one at a time and
DuckDB 7~ms. Polars sits near DuckDB, slower on the grouped shapes; pandas pays for scanning the whole in-memory table with a
boolean mask on every query and is the slowest on four shapes of five. The hand-written numpy engine is the fastest on the
range, the bars and the cross-section, and close to the fastest on the other two: knowing the layout (sorted, with an offset per
symbol) beats any general engine, at the cost of writing and maintaining it for every query.

\omcode{code/firm/tsbench/firm_tsbench.py}{255}{274}{Cold and warm timings: a fresh engine's first run of each query is the
cold one; warm runs are interleaved, every engine once per query per round.}\label{lst:pl:time-series-databases-and-the-alternatives:bench}

\begin{example}[The query mix decides]\label{ex:pl:time-series-databases-and-the-alternatives:mix}
A workload of point lookups and cross-sections, a fraction $f$ of them cross-sections, costs $(1-f)\,t_{\mathrm{point}} +
f\,t_{\mathrm{xsec}}$ per query on each engine. With the medians above, DuckDB and SQLite cost the same at $f = 1.2\%$: as soon as
more than about one query in a hundred is a cross-section, the columnar engine is faster on the whole workload, although it
loses nearly eighty to one on the lookups.
\end{example}

\section{A survey of the field}

\begin{dated}{2026-09}{Time-series and analytical engines in use for market data}\label{dat:pl:time-series-databases-and-the-alternatives:survey}
As described by their own documentation and licence files (consulted September 2026): \textbf{kdb+} (KX) is ``a high-performance
cross-platform historical timeseries columnar database'', an in-memory compute engine and a realtime streaming processor with its
own language, q; all use requires a licence (commercial, or a free community edition of its successor, KDB-X, with usage
restrictions). \textbf{ClickHouse} is an open-source column-oriented database server for real-time analytical reports (Apache
2.0). \textbf{QuestDB} is an open-source \omterm{def:pl:time-series-databases-and-the-alternatives:tsdb}{time-series database} with SQL over memory-mapped, time-partitioned columns and Parquet
(Apache 2.0). \textbf{TimescaleDB} is a PostgreSQL extension for time-series analytics with a columnstore (Apache 2.0 outside
parts under the Timescale Licence). \textbf{InfluxDB~3} is an open-source \omterm{def:pl:time-series-databases-and-the-alternatives:tsdb}{time-series database} built on Arrow, DataFusion and
Parquet (MIT or Apache 2.0). \textbf{ArcticDB} (Man Group) is a serverless dataframe database for Python storing to object
storage or LMDB (Business Source Licence 1.1, each version converting to Apache 2.0 after at most four years). \textbf{DuckDB}
and \textbf{Polars} are embedded, under the MIT licence.
\end{dated}

\begin{table}[ht]
\centering\small
\begin{tabular}{lp{2.7cm}p{3.3cm}p{4.3cm}}
\toprule
engine & deployment & storage & strength on market data \\
\midrule
kdb+ & server, in-memory & columnar, time-partitioned & tick capture and queries in one system \\
ClickHouse & server & columnar & large scans and aggregations \\
QuestDB & server & columnar, time-partitioned & ingestion and time-range SQL \\
TimescaleDB & PostgreSQL extension & rows and a columnstore & relational data beside time series \\
InfluxDB 3 & server & Arrow and Parquet & monitoring-style metrics \\
ArcticDB & library, object store & versioned dataframes & dataframe research history \\
DuckDB, Polars & library & files (Parquet, Arrow) & research on a \omterm{def:pl:a-tick-store-on-open-formats:store}{tick store} \\
SQLite & library & rows, B-tree & small reference tables, lookups \\
\bottomrule
\end{tabular}
\caption{The engines of the dated box, by deployment and storage, and the part of the market-data workload each is built for;
the last column is this book's reading of their designs, not a benchmark.}\label{tab:pl:time-series-databases-and-the-alternatives:survey}
\end{table}

The choice follows from the query mix and from who writes. A firm whose market-data queries are research scans over history,
written once a day by a compaction job, is served by files and an embedded engine; a firm that needs many concurrent writers and
readers of live data with access control buys or builds a server; a firm whose researchers live in dataframes and want versioned
datasets looks at a dataframe store. Whatever the choice, the benchmark of this chapter, on the firm's own shapes, decides between
candidates.

\section{Tutorial: five engines, one answer}\label{tut:pl:time-series-databases-and-the-alternatives:tut}

\begin{tutorial}
\textbf{Goal.} Load one dataset into five engines, prove that they agree, and time them properly. \textbf{End state:}
\cref{fig:pl:time-series-databases-and-the-alternatives:times} and the break-even mix of
\cref{ex:pl:time-series-databases-and-the-alternatives:mix}.
\begin{tutsteps}
\item \textbf{Data}: \texttt{pl\_tsbench.dataset()}: the chapter-3 month of quotes and 40\,000 trades; \texttt{prepare()} writes
them as Parquet partitioned by date and sorted by symbol and time.
\item \textbf{Agree}: \texttt{answers()} builds the five engines and checks the five shapes; every check must be true.
\item \textbf{Time}: \texttt{bench\_tsbench.py} runs \texttt{firm\_tsbench.bench}
(\cref{lst:pl:time-series-databases-and-the-alternatives:bench}): cold runs, then seven interleaved rounds.
\item \textbf{Decide}: compute the break-even share of cross-sections between the fastest lookup engine and the fastest scan
engine.
\end{tutsteps}
\textbf{What to change next.} Drop SQLite's index and rerun the point lookup; add a sixth shape (the volume-weighted price of
each symbol's trades on one day) to every engine and check they agree before timing it.
\end{tutorial}

\section{Build: the engine benchmark}\label{bld:pl:time-series-databases-and-the-alternatives:tsbench}

\begin{build}
\bfield{Purpose} The firm's standing benchmark for any engine proposed for market data: same data, same shapes, answers checked,
timings interleaved.
\bfield{Interface} \texttt{QUERIES}, \texttt{ENGINES}, \texttt{Params(symbol, date, t0, t1, point\_t)};
\texttt{write\_dataset(quotes, trades, root)}; \texttt{make\_engine(name, quotes, trades, root)} with \texttt{run(query, params)}
returning sorted, rounded rows; \texttt{check(engines, params)}; \texttt{bench(factory, params, rounds)} returning cold, median
and spread per engine and query.
\bfield{Rules} No timing before every engine agrees on every shape; rounds interleave engines; cold runs are fresh instances;
results are compared after sorting and rounding to six decimals, never by eye.
\bfield{Acceptance tests} \texttt{code/firm/tsbench/tests/}: all five engines agree on every shape on a small dataset; a planted
wrong engine is caught by the check; the as-of join returns the quote one nanosecond before each trade; the benchmark returns a
cold and a warm time for every engine and shape.
\bfield{Stretch} A server engine behind the same interface; memory as well as time; the benchmark at several data sizes, to
estimate each engine's $c_0$ and $c_r$.
\end{build}

\begin{omsources}
\item P.~Boncz, M.~Zukowski and N.~Nes, ``MonetDB/X100: hyper-pipelining query execution'', \emph{CIDR}, 2005.
\item D.~J.~Abadi, D.~S.~Myers, D.~J.~DeWitt and S.~R.~Madden, ``Materialization strategies in a column-oriented DBMS'',
\emph{ICDE}, 2007.
\item M.~Raasveldt and H.~Mühleisen, ``DuckDB: an embeddable analytical database'', \emph{SIGMOD}, 2019.
\item The documentation and licence files of the engines in the dated box.
\end{omsources}

\section{Exercises}

\begin{exercise}[$\star$]\label{exo:pl:time-series-databases-and-the-alternatives:1}
From \cref{fig:pl:time-series-databases-and-the-alternatives:times}, by what factor is SQLite faster than DuckDB on the point
lookup, and by what factor slower on the cross-section?
\end{exercise}

\begin{exercise}[$\star$]\label{exo:pl:time-series-databases-and-the-alternatives:2}
Why does an index on (symbol, time) make the point lookup almost free, and why does it not help the cross-section?
\end{exercise}

\begin{exercise}[$\star$]\label{exo:pl:time-series-databases-and-the-alternatives:3}
A session of 6.5 hours has how many one-minute buckets? The bars query returned 359 rows for S007's day: what happened to the
others?
\end{exercise}

\begin{exercise}[$\star\star$]\label{exo:pl:time-series-databases-and-the-alternatives:4}
Why are the engines interleaved within each round rather than each run seven times in a row?
\end{exercise}

\begin{exercise}[$\star\star$]\label{exo:pl:time-series-databases-and-the-alternatives:5}
Explain, with \cref{prop:pl:time-series-databases-and-the-alternatives:shapes}, why DuckDB's point lookup (3.8~ms) is slower than
its range query (2.2~ms), although the range returns 153 rows.
\end{exercise}

\begin{exercise}[$\star\star$]\label{exo:pl:time-series-databases-and-the-alternatives:6}
Derive the break-even share of cross-sections of \cref{ex:pl:time-series-databases-and-the-alternatives:mix} from the measured
medians.
\end{exercise}

\begin{exercise}[$\star\star\star$]\label{exo:pl:time-series-databases-and-the-alternatives:7}
\emph{Coding.} Check all five engines on symbol S023, day 12 (the same windows, shifted to that day). Do they agree, and how many
rows does each shape return?
\end{exercise}

\begin{exercise}[$\star\star\star$]\label{exo:pl:time-series-databases-and-the-alternatives:8}
\emph{Find the flaw.} ``We compared the vendor's \omterm{def:pl:time-series-databases-and-the-alternatives:tsdb}{time-series database} with DuckDB on the vendor's demonstration query: one run of
each, ours first. The vendor's was three times faster, so we are buying it.''
\end{exercise}

\section{Problem: The Right Tool for Each Question}

\begin{problem}[{Weekend problem --- choosing an engine for the research platform}]\label{pb:pl:time-series-databases-and-the-alternatives:1}
The chapter's month of quotes and trades, the five engines of \texttt{firm.tsbench} and the measured medians of
\cref{fig:pl:time-series-databases-and-the-alternatives:times}.

\medskip\noindent\textbf{Part I --- Shapes.}
\begin{enumerate}
\item Name the five query shapes and give a market-data question of each.
\item Which shapes touch few rows, and which touch many?
\item What layout serves the range and the bars shapes?
\item What does a \omterm{def:pl:time-series-databases-and-the-alternatives:tsdb}{time-series database} add to the \omterm{def:pl:a-tick-store-on-open-formats:store}{tick store} of chapter~4?
\item What is \omterm{def:pl:time-series-databases-and-the-alternatives:bucket}{downsampling}, and where does it belong in a \omterm{def:pl:capturing-and-storing-tick-data:tier}{retention policy}?
\end{enumerate}

\medskip\noindent\textbf{Part II --- Engines.}
\begin{enumerate}[resume]
\item What distinguishes an embedded engine from a server, and what does each give up?
\item What is vectorised execution, and why does it make scans fast?
\item What is \omterm{def:pl:time-series-databases-and-the-alternatives:vectorised}{late materialisation}?
\item Why is a B-tree lookup faster than a vectorised scan of a sorted file for one row?
\item Which engine is fastest on each shape here, and which is slowest?
\end{enumerate}

\medskip\noindent\textbf{Part III --- Measurement.}
\begin{enumerate}[resume]
\item How many rows does each shape return, and why check them before timing?
\item What are SQLite's and DuckDB's medians on the point lookup and the cross-section?
\item Why is the hand-written numpy engine fastest on four shapes, and why is it not the answer?
\item What does the cold run measure that the warm median does not?
\item What must be stated with the numbers for them to mean anything?
\end{enumerate}

\medskip\noindent\textbf{Part IV --- The verdict.}
\begin{enumerate}[resume]
\item State the \emph{named result}: the fastest engine for each shape, and the query mix at which the columnar engine beats the
row store.
\item A research team's workload is 95\% cross-sections and bars, 5\% point lookups: which engine?
\item A service answering last-quote lookups for an order router: which engine, and why not the research one?
\item When is a \omterm{def:pl:time-series-databases-and-the-alternatives:tsdb}{time-series database} server the right purchase?
\item In one sentence: what decides which engine is fastest?
\end{enumerate}
\end{problem}

\section{Interview questions}

\begin{interviewq}[$\star$ \iqroles{developer, researcher}]\label{iq:pl:time-series-databases-and-the-alternatives:1}
What is a \omterm{def:pl:time-series-databases-and-the-alternatives:tsdb}{time-series database}, and do you need one for \omterm{def:pl:capturing-and-storing-tick-data:tick}{tick data}?
\end{interviewq}

\begin{interviewq}[$\star\star$ \iqroles{developer}]\label{iq:pl:time-series-databases-and-the-alternatives:2}
Why is a columnar engine fast at aggregations and slow at single-row lookups?
\end{interviewq}

\begin{interviewq}[$\star\star$ \iqroles{developer}]\label{iq:pl:time-series-databases-and-the-alternatives:3}
How would you benchmark two databases for your market data so that the result means something?
\end{interviewq}

\begin{interviewq}[$\star\star$ \iqroles{developer, researcher}]\label{iq:pl:time-series-databases-and-the-alternatives:4}
When would you use an embedded engine such as DuckDB rather than a database server?
\end{interviewq}

\begin{interviewq}[$\star\star\star$ \iqroles{developer}]\label{iq:pl:time-series-databases-and-the-alternatives:5}
Your firm is asked to replace its proprietary tick database. How do you decide what to replace it with?
\end{interviewq}
